\section{Physical transport of the common parent ground space}%
\label{sec:peps_region_physical_ground_space_transport}

\begin{definition}[Product physical map on the full graph]%
    \label{def:peps_global_physical_map}
    \lean{TNLean.PEPS.fullRegionConfigEquiv,
        TNLean.PEPS.fullRegionPhysicalEquiv,TNLean.PEPS.globalPhysicalMap,
        TNLean.PEPS.globalPhysicalMap_apply,TNLean.PEPS.globalPhysicalEquiv,
        TNLean.PEPS.globalPhysicalEquiv_apply}
    \leanok
    \uses{def:peps_physical_deformation}
    For physical maps $F_v:\C^d\to\C^e$ at the vertices of a finite
    graph, their product map $F_V=\bigotimes_{v\in V}F_v$ acts by
    \begin{align}
        (F_V\psi)(\tau)=\sum_{\sigma:V\to\{0,\ldots,d-1\}}
        \left(\prod_{v\in V}(F_v)_{\tau_v\sigma_v}\right)\psi(\sigma).
    \end{align}
    If each $F_v$ is invertible, then $F_V$ is a linear isomorphism.
\end{definition}

\begin{theorem}[Physical deformation of the closed PEPS vector]%
    \label{thm:peps_physical_deformation_closed_state}
    \lean{TNLean.PEPS.stateCoeff_physicalDeform}
    \leanok
    \uses{def:peps_global_physical_map,def:peps_physical_deformation}
    The closed PEPS vectors satisfy
    $\ket{\Psi(A^F)}=F_V\ket{\Psi(A)}$.
\end{theorem}

\begin{proof}\leanok
    Expand both coefficients as finite sums over virtual labels and
    physical configurations. Interchange the two sums and distribute
    the product of physical coefficients at the vertices. This gives
    the defining local physical sums of $A^F$.
\end{proof}

\begin{theorem}[Slices under a product physical map]%
    \label{thm:peps_global_physical_map_slice}
    \lean{TNLean.PEPS.physicalProduct_assembleRegion,
        TNLean.PEPS.regionSliceMap_globalPhysicalMap}
    \leanok
    \uses{def:peps_global_physical_map}
    Fix a region $R\subseteq V$ and write $R^c=V\setminus R$.
    For a complementary physical configuration $\beta$, let
    $\psi_\beta(\alpha)=\psi(\alpha,\beta)$.
    Every complementary output configuration $\tau$ satisfies
    \begin{align}
        (F_V\psi)_\tau
        =\sum_\beta\left(\prod_{v\in R^c}(F_v)_{\tau_v\beta_v}\right)
        F_R\psi_\beta.
    \end{align}
\end{theorem}

\begin{proof}\leanok
    Split each input configuration into its restrictions to $R$ and
    $R^c$. The coefficient product factors over these two sets.
    Interchanging the two finite sums gives the stated identity.
\end{proof}

\begin{theorem}[Transport of the common parent ground space]%
    \label{thm:peps_region_parent_ground_space_physical_transport}
    \lean{TNLean.PEPS.mem_regionParentGroundSpace_iff,
        TNLean.PEPS.globalPhysicalMap_mem_regionParentGroundSpace,
        TNLean.PEPS.globalPhysicalEquiv_mem_regionParentGroundSpace_iff,
        TNLean.PEPS.map_regionParentGroundSpace_physicalDeform,
        TNLean.PEPS.regionParentGroundSpacePhysicalEquiv,
        TNLean.PEPS.finrank_regionParentGroundSpace_physicalDeform,
        TNLean.PEPS.map_ker_regionParentHamiltonian_physicalDeform}
    \leanok
    \uses{def:peps_global_physical_map,def:peps_region_parent_hamiltonian}
    Let $\{R_i\}_{i\in I}$ be any family of regions, and let
    $\mathcal G(A)$ be the space of full physical vectors whose slices
    on every $R_i$ belong to $\mathcal G_{R_i}(A)$.
    If each $F_v$ is invertible, then
    $\mathcal G(A^F)=F_V\mathcal G(A)$.
    Restriction of $F_V$ is a linear isomorphism between these common
    ground spaces, so their dimensions agree.
    For a finite family of positive parent interactions $h_i$,
    define $h_i^F=(F_{R_i}^{-1})^\dagger h_iF_{R_i}^{-1}$.
    The corresponding full parent Hamiltonians satisfy
    $\ker H^F=F_V\ker H$.
\end{theorem}

\begin{proof}\leanok
    \uses{thm:peps_global_physical_map_slice,
        thm:peps_region_physical_deformation,
        thm:peps_deformed_region_interaction,
        thm:peps_region_parent_common_kernel}
    If every slice $\psi_\beta$ lies in $\mathcal G_{R_i}(A)$,
    then $F_{R_i}\psi_\beta$ lies in
    $\mathcal G_{R_i}(A^F)$. The slice identity and closure under
    finite linear combinations give $F_V\mathcal G(A)\subseteq
        \mathcal G(A^F)$. Applying the inverse physical maps gives the
    reverse inclusion. Restriction of $F_V$ is therefore a linear
    isomorphism. The full Hamiltonian kernel is the common regional
    ground space for every choice of positive interactions with the
    specified local kernels, giving the final identity.
\end{proof}
